\section{The relation index of \citet{ren2024semantic}}
\label{app:ri}

This appendix restates the SIH criterion as defined in the original paper (their
\S3, Eq.~2--3), so that the main text can stay short, and lists the evidence behind
the limitations noted in \S\ref{sec:sih}. Our own re-implementation, with its
controls and calibration, is described in \S\ref{sec:method-ri}.

\paragraph{Notation.}
The model has $L$ layers of $H$ heads. For head $h$, $W^h_{OV}=W^h_v W^h_o$ is the OV
matrix, $W^h_{QK}=W^h_q (W^h_k)^{\top}$ the QK matrix, $W_e$ the input embedding and
$W_U$ the unembedding. For an input $[t_1,\ldots,t_n]$, $x_i=t_i W_e$ is the
\emph{input embedding} of token $t_i$. A relation is a triplet $T=(t_s,r,t_o)$ with
head token at position $s$ and tail token at position $o$.

\paragraph{Step 1: attention condition (QK).}
At a current position $j\ge\max(s,o)$, let $A^h_j$ be head $h$'s attention over
positions $1..j$, written in the original paper as $\mathrm{softmax}(x_j W^h_{QK}
x^{\top}/\sqrt{d_h})$ in the notation of Eq.~1 there. The head qualifies for $(T,j)$
if
\begin{equation*}
s=\arg\max_{k\le j} A^h_{j,k}
\quad\text{and}\quad
\frac{A^h_{j,s}}{\max_{k\ne s} A^h_{j,k}}>\tau,\qquad \tau=2.2 .
\end{equation*}
The paper's formula writes the attention in terms of the embedding matrix $x$; in our
implementation $A^h_j$ is the attention pattern actually produced by the model's
forward pass at the head's layer (captured from the attention kernel), and only the
OV projection of Step~2 uses the raw embedding.

\paragraph{Step 2: output condition (OV).}
The head's output at $j$ is projected onto the vocabulary from the input embedding of
the current token,
\begin{equation*}
p^{h,j}=\mathrm{softmax}\!\left(x_j W^h_{OV} W_U\right)\in\mathbb{R}^{|V|}.
\end{equation*}
Only the tokens $t_k$, $k\le j$, that appear in the context are kept (duplicates
counted once). Their probabilities are centred and clipped,
$q^{h,j}_{t_k}=\max\!\big(0,\;p^{h,j}_{t_k}-\mathbb{E}_{k'\le j}[p^{h,j}_{t_{k'}}]\big)$,
and the score for $(T,j)$ is the share of the clipped mass on the tail token,
\begin{equation*}
a^{h,j}_T=\frac{q^{h,j}_{t_o}}{\sum_{k\le j} q^{h,j}_{t_k}} .
\end{equation*}

\paragraph{Worked example.}
Table~\ref{tab:ri-example} carries the computation through for one event, with
illustrative numbers. The context is ``Anna is the mother of Ben .'' and the current
token is the period ($j{=}7$), so seven distinct tokens are visible. Suppose head $h$
attends most to \emph{Anna} with a margin above $\tau$, so the event qualifies with
$t_s{=}$Anna. The raw-embedding projection of the period is pushed through $W^h_{OV}W_U$
and the softmax is restricted to the visible tokens; the mean of the seven values,
$\bar p{=}0.143$, is subtracted and negatives are clipped. The tail token \emph{Ben}
holds $0.52$ of the clipped mass, which is the event score. Two properties of the
score are visible here: it depends only on the current token's embedding and on
which tokens are visible, so every event of this head at a period with the same
visible set receives the same vector $p$; and the tail competes with all visible
tokens, including \emph{is} and \emph{the}, not with other names.

\begin{table}[h]
\centering\scriptsize
\setlength{\tabcolsep}{2pt}
\begin{tabular}{@{}lrrrrrrr@{}}
\toprule
$t_k$ & Anna & is & the & mother & of & Ben & . \\
\midrule
$p^{h,j}_{t_k}$ & 0.30 & 0.05 & 0.05 & 0.15 & 0.05 & 0.32 & 0.08 \\
$q^{h,j}_{t_k}$ & 0.157 & 0 & 0 & 0.007 & 0 & 0.177 & 0 \\
\bottomrule
\end{tabular}
\caption{Illustrative event score. $q=\max(0,p-\bar p)$ with $\bar p=0.143$;
$\sum_k q_k=0.341$; $a^{h,j}_T=q_{\text{Ben}}/\sum_k q_k=0.177/0.341=0.52$. The numbers are chosen for
the example, not measured.}
\label{tab:ri-example}
\end{table}

\paragraph{Step 3: aggregation.}
$\mathrm{RI}_h$ is the mean of $a^{h,j}_T$ over all qualifying $(T,j)$ for a given
relation type. There is no threshold that turns an index into a label: the paper
reports heatmaps of $\mathrm{RI}_h$ per relation and calls heads with salient values
SIHs. An appendix of the original paper shows that $\tau\in\{2.0,2.5\}$ yields a
similar set of salient heads.

\paragraph{Setting of the original study.}
InternLM2-1.8B (24 layers, 16 heads); the AGENDA test set (1,000 abstracts with
annotated knowledge graphs); three syntactic dependencies (subject, object,
modifier) and seven relation types; multi-token entities replaced by single capital
letters and frequent function words removed before scoring.

\paragraph{Evidence for the limitations in \S\ref{sec:sih}.}
\begin{itemize}\setlength\itemsep{1pt}
\item[(i)] The word ``ablation'' occurs once in the original paper and refers to
varying $\tau$ (their Appendix~A). No activation, weight or attention intervention is
reported.
\item[(ii)] The training analysis (their \S4.2) uses a model the authors train
themselves ($d{=}2048$, 40k steps on SlimPajama, checkpoints every 200 steps), not
the released InternLM2-1.8B in which the SIHs of \S3 are identified. In that analysis
the margin condition on $\tau$ is dropped (``it is too challenging to find attention
heads having an extremely high attention probability''), and the plotted heads are
``sampled \ldots with an increasing relation index''. The inference is stated as
``we infer that the formation of semantic induction heads plays a crucial role in the
development of the ICL ability.''
\item[(iii)] The projection $x_j W^h_{OV} W_U$ uses $x_j=t_j W_e$. The paper justifies
this by quoting \citet{elhage2021mathematical}: ``both the QK circuit and OV circuit
are directly performed on input embeddings.'' That statement describes attention-only
models without MLPs and normalization, where the residual at every layer is a sum of
head outputs added to the embedding; in a 24- or 32-layer model with MLPs and
layer normalization the value input of head $h$ at position $j$ is the normalized
residual $\mathrm{LN}(r^{\ell}_j)$, and the head's output at $j$ is
$\sum_k A^h_{j,k}\,\mathrm{LN}(r^{\ell}_k)W^h_{OV}$, a function of the attended
positions $k$.
\item[(iv)] In $a^{h,j}_T$ the tail competes with every earlier distinct token, not
with tokens of the same type (other entities), and no reference distribution for
$\mathrm{RI}_h$ is reported.
\end{itemize}
